% Generated draft: requires mathematical and source review.
\documentclass[11pt]{article}
\usepackage[margin=25mm]{geometry}
\usepackage{fontspec}
\setmainfont{Latin Modern Roman}
\usepackage{amsmath,amssymb,mathtools,unicode-math}
\setmathfont{Latin Modern Math}
\usepackage{enumitem,needspace,xurl,hyperref}
\setlength{\parindent}{0pt}
\setlength{\parskip}{5pt}
\setlength{\emergencystretch}{3em}
\newcommand{\N}{\mathbb{N}}
\newcommand{\Z}{\mathbb{Z}}
\newcommand{\Q}{\mathbb{Q}}
\newcommand{\R}{\mathbb{R}}
\newcommand{\C}{\mathbb{C}}
\newcommand{\F}{\mathbb{F}}
\newcommand{\A}{\mathbb{A}}
\newcommand{\PP}{\mathbb{P}}
\newcommand{\E}{\mathbb{E}}
\newcommand{\eps}{\varepsilon}
\newcommand{\dd}{\,\mathrm{d}}
\begin{document}
\section*{UnsolvedMath ID: 30005323}
\textbf{Source problem code:} catalogue entry
\subsection{Definitions and mathematical statement}
For a nonnegative integer \(n\), a partition \(\lambda\vdash n\) is a finite nonincreasing sequence of positive integers with sum \(n\); the empty partition is allowed when \(n=0\). Let \(S_n\) denote the set of permutations of \(\{1,\ldots,n\}\), with \(S_0\) consisting of the empty permutation, so that \(|S_n|=n!\). These conventions explain the notation; the supplied record does not specify the intended range of \(n\) or whether \(n=0\) is included.

The supplied problem requests a bijective proof of
\[
\sum_{\lambda\vdash n} f_\lambda^{\,2}=n!
\]
for the intended values of \(n\). However, the record does not define \(f_\lambda\) or identify the objects it counts. Consequently, the intended bijection cannot be specified completely from this record.

More precisely, if the missing source definitions provide finite sets \(\mathcal F_\lambda\) with \(|\mathcal F_\lambda|=f_\lambda\), a bijective formulation of the displayed identity would request, for every \(n\) in the intended range, an explicit bijection
\[
\bigsqcup_{\lambda\vdash n}
\bigl(\mathcal F_\lambda\times\mathcal F_\lambda\bigr)
\longrightarrow S_n.
\]
This is a conditional formulation: the sets \(\mathcal F_\lambda\), the range of \(n\), and any requirements inherited from the Hurwitz-number setting remain unspecified.

The original excerpt introduces the identity by asserting \(H_{n,0}=1\) and saying that the case \(\ell=0\) reduces to this bijective task. Neither \(H_{n,\ell}\) nor the role and domain of \(\ell\) are defined in the supplied record. These qualifications must be recovered before identifying this draft with the full source problem.
\subsection{Short English statement}
The record asks for a bijective explanation of \(\sum_{\lambda\vdash n}f_\lambda^2=n!\), presented as the \(\ell=0\) case of a Hurwitz-number problem with \(H_{n,0}=1\). The counted objects, parameter range, and surrounding Hurwitz-number definitions are missing, so the intended problem remains only partially specified.
\subsection{Sources}
\begin{itemize}
\item Supplied local UnsolvedMath registry record, permanent ID \(30005323\), source problem code \texttt{catalogue entry}, titled ``Bijective Proof of a Hurwitz Number Identity.'' The original excerpt is labeled ``Open problem 1.'' No individual proposer is supplied; the proposed year is recorded as 2022.
\item Supplied report citation: \emph{Enumerative Combinatorics (2023)}, Oberwolfach Reports, pp.~3215--3305, DOI \texttt{10.4171/owr/2022/57}. Source URL and literature-source link labeled ``original OWR report'': \url{https://doi.org/10.4171/owr/2022/57}. This is the supplied authoritative source reference; its contents have not been independently consulted here.
\item Literature assessment dated exactly \texttt{2026-08-22}: the supplied curation labels the problem open and reports that no authoritative source located by its cutoff proves or disproves the specific extracted claim. This is supplied dated curation, not independently verified mathematical or literature status. The background explicitly requests independent verification of time-sensitive status and cited claims.
\item The registry labels the statement \texttt{corrected\_verified} and states that only formula delimiters, notation, grammar, and typography were repaired, with mathematical content and scope preserved. This supplied editorial qualification does not resolve the missing definitions or certify the formulation below.
\end{itemize}
\end{document}
